\section{Introduction}
An advisory can describe an authorization bypass or unsafe operation without identifying the responsible source lines. Locating those lines requires navigating a repository, often across an externally reachable entry point and a security-critical operation. A language model can use advisory clues, repository evidence, and prior knowledge to propose those locations. This acquisition can be task-independent, driven by a retrieval query, or adapted through successive tool interactions.

The comparison is more involved than assigning each method the same context limit. A file may be eligible for a method but never selected; a selected span may contain little relevant evidence; and an informative interaction may still end without a usable answer. Replaying an interactive history also consumes provider input without delivering new repository information. Consequently, localization recall alone cannot establish either acquisition effectiveness or resource efficiency.

We examine these distinctions through three implemented configurations for advisory-guided localization: a read-only exploration Agent, single-round hybrid retrieval (RAG), and lexicographic context packing (Pack). All use DeepSeek V4 Pro and a 50,000-token ceiling on the repository-information ledger. Inputs combine a VulnGym entry title with a truncated public advisory description. Entry titles can include filenames appended during dataset annotation; 47 of the 80 supplied titles contain an EP/CO target basename. We therefore study localization with these dataset-supplied titles, rather than assuming an unmodified advisory input. Explicit entry-point (EP) and critical-operation (CO) annotation fields are supplied to the evaluator, while their filename clues may also occur in input text~\cite{vulngym,vulngymdata}. The model predicts up to five file/line intervals; the task concerns known vulnerabilities.

Three research questions organize the study:
\begin{itemize}
\item \textbf{RQ1:} How do the complete configurations differ in file and line recall when every scheduled model outcome is retained?
\item \textbf{RQ2:} How do source delivery, EP/CO annotation coverage, prediction extent, and output validity co-occur?
\item \textbf{RQ3:} What delivered information, cumulative provider input, reported fees, and online execution time accompany these outcomes?
\end{itemize}

The evidence consists of an original study (S0) and a supplementary study (S1), each comprising 720 outcomes: 80 advisory-distinct entries, three configurations, and three repeats. S1 reuses the known sample after S0 inspection and jointly revises source eligibility, indexing, code presentation, and the Agent interface. We preserve both studies and report them separately. S1 supplies a comparison with shared nonblank source eligibility; it is neither an independent-sample confirmation nor an ablation of a single change.

Our contribution is an empirical characterization of how the interpretation of localization scores changes when source availability, annotation roles, interval width, and output validity are examined jointly. In S1, most Agent annotations left uncovered despite source delivery are entry points, a role not explicitly requested separately in the task. Most RAG hits without annotated source depend on broad predictions. Schema acceptance also differs from repository-valid locations. These observations identify which distinctions an aggregate recall score leaves unresolved.

We provide run-linked manifests, delivered spans, model messages, retained failures, and descriptive analysis scripts. A complete-denominator decomposition reports how output failures contribute to the Agent--RAG file contrast. The evidence supports evaluation of the measured title-assisted task: triaging locations for a supplied known vulnerability with filename clues potentially available. Performance in hint-free discovery and the causes of individual model decisions are separate questions.
